%%%%%%%%%%  scan idx 232  |  folio 219  %%%%%%%%%%
\SGAPrintPage{219}%
$\operatorname{Ell}^{\underline n}(F)$ of
$\operatorname{Comp.diff}^{\underline n}(F)$ formed by the elliptic complexes of order $\underline n$ is open; it is this space (which in general is not locally compact) that plays the role of a space of universal local parameters.

\sgaStableTarget{sga6.e02.para.3-2-4}3.2.4. Nevertheless, even should the generalization of (3.2.1) mentioned above prove impossible, we can always associate to $E^\bullet$, under the sole hypothesis that the fibration $p$ is proper and locally trivial, an element of $K^\bullet(S)$ deserving the name analytic index of $E^\bullet$. Indeed, suppose (for simplicity\SGAAuthorFootnote{1}{for the general case, use the notion of a \textquotedblleft weak Hilbert bundle\textquotedblright{} or continuous field of Hilbert spaces from [12].}) that we are in the \textquotedblleft local\textquotedblright{} situation where $X=S\times V$ ($V$, a compact variety), with $E^\bullet$ given by a continuous map from $S$ to $\operatorname{Ell}^{\underline n}(F)$ (3.2.3). Assertions (a') and (b) in the proof of (3.2.1) are in fact valid without any hypothesis on $S$: it is therefore reasonable to put, with the notation of (loc.cit.):
\[
  \operatorname{ind.an}(E^\bullet)=C\ell(\underline F_r^\bullet)\in K^\bullet(S),
\]
this element being independent of $r$.

\sgaStableTarget{sga6.e02.para.3-2-5}3.2.5. Suppose $S$ is paracompact, and denote by $H$ a complex Hilbert space of infinite Hilbert dimension, and by $I(H)$ the subspace of $L(H,H)$ formed by the operators \textquotedblleft with index\textquotedblright{}. Then the complexes $F_r^\bullet$ define [12] a canonical element in the group $[S,I(H)]$ of homotopy classes of continuous maps from $S$ to $I(H)$, but the writer does not know, for noncompact $S$ (even when $S$ is locally compact), the relations between the groups $[S,I(H)]$ and $K^\bullet(S)$ (is there a natural arrow from one to the other? are they isomorphic?), and therefore a fortiori between the two indices thus defined.

For a generalization, under the hypotheses of (3.2.1), of the Atiyah--Singer index theorem, we refer to the recent work of Shih [17], or to work in preparation by M. F. Atiyah and I. M. Singer.
%%%%%%%%%%  scan idx 233  |  folio 220  %%%%%%%%%%
\SGAPrintPage{220}%
\sgaStableTarget{sga6.e02.bibliography}\subsection*{BIBLIOGRAPHY}
\begin{enumerate}[label={[\arabic*]}]
\item
Artin, M., Grothendieck, A., and Verdier, J-L., Cohomologie étale des schémas, to appear in North Holl. Pub. Comp., (cited as SGA 4).
\item Atiyah, M. and Segal, G., secret papers.
\item Cartan, H. Funzioni e varietà complesse, Faisceaux analytiques cohérents, Course at the Centro Internazionale Matematico Estivo, 1964, Roma.
\item Cartan, H., and Schwartz, L., Théorème d'Atiyah--Singer, Seminar 1963-64, Institut Henri Poincaré, Paris.
\item Frisch, J., Fonctions analytiques sur un ensemble semi-analytique, note in CRAS Paris, t. 260, p. 2974-2976 (1965).
\item Godement, R., Théorie des Faisceaux, Hermann (1958), (cited as T.F).
\item Grothendieck, A., Cohomologie locale des faisceaux cohérents et théorèmes de Lefschetz locaux et globaux, Séminaire de Géométrie Algébrique de l'IHES 1962, to appear in North Holl. Pub. Comp., (cited as SGA 2).
\item Grothendieck, A., Fondements de la Géométrie Analytique, Séminaire Cartan 1959-60, Institut Henri Poincaré, Paris.
\item Hartshorne, R., Residues and Duality, Lecture Notes, Springer n\textsuperscript{o}~20, cited as [H].
\item Houzel, C., Géométrie Analytique Locale, Séminaire Cartan 1959-60, Institut Henri Poincaré, Paris.
\item Illusie, L., Complexes de fibrés et définitions du foncteur $K$, Séminaire Shih Weishu, IHES 1964-65.
\item Illusie, L., Complexes quasi-acycliques directs de fibrés banachiques, note in CRAS Paris, t. 260, p. 6499-6502.
\item Kiehl, R., Theorem A und Theorem B in der nichtarchimedischen Funktionentheorie, Invent. math. 2, p. 256-273 (1967).
%%%%%%%%%%  scan idx 234  |  folio 221  %%%%%%%%%%
\SGAPrintPage{221}%
\item
Lang, S., Introduction to differentiable manifolds,
\item \L{}ojasiewicz, S., Ann. della Scuola Norm. Sup. Pisa, serie III, 18, fasc. IV, 1964.
\item Palais, R., Lectures at the Princeton Seminar 1963-64 on the index theorem.
\item Shih, W., Fiber cobordism and the index of a family of elliptic differential operators,
\item Verdier, J.-L., Les catégories dérivées des catégories abéliennes, to appear in North Holl. Pub. Comp.
\end{enumerate}

% Original English continuation definitions, retained in their original order.
\providecommand{\Ob}{\operatorname{Ob}}
\providecommand{\qcoh}{\mathrm{qcoh}}
\providecommand{\coh}{\mathrm{coh}}
\providecommand{\torf}{\mathrm{torf}}
\providecommand{\toramp}{\operatorname{tor.amp}}
\providecommand{\amp}{\operatorname{amp}}
\providecommand{\Id}{\operatorname{Id}}
\providecommand{\pr}{\operatorname{pr}}
\providecommand{\coker}{\operatorname{coker}}
\providecommand{\im}{\operatorname{Im}}
\providecommand{\Spec}{\operatorname{Spec}}

% Reviewed French continuation definitions, retained verbatim.
% Exposé III additions
\renewcommand{\Ob}{\operatorname{ob}}
\providecommand{\qcoh}{\mathrm{qcoh}}
\providecommand{\coh}{\mathrm{coh}}
\providecommand{\torf}{\mathrm{torf}}
\providecommand{\toramp}{\operatorname{tor.amp}}
\providecommand{\amp}{\operatorname{amp}}
\providecommand{\Id}{\operatorname{Id}}
\providecommand{\pr}{\operatorname{pr}}
\providecommand{\coker}{\operatorname{coker}}
\providecommand{\im}{\operatorname{Im}}
\providecommand{\Spec}{\operatorname{Spec}}

